\chapter{Kolmogorov, Cover and the Information in a Solution}
\label{ch:kolmogorov}

Shannon measured the information in a source; Kolmogorov\index{Kolmogorov, Andrei} measured the information in a single object. The difference matters for this book because a queens solution is a single object, and the question whether it can be found, counted or sampled efficiently is, in Kolmogorov's terms, the question how many bits it takes to describe one and how long it takes to decode the description. Tom Cover\index{Cover, Thomas} spent a career on the bridge between the two measures, and three of his results, enumerative coding, universal gambling and the account of Kolmogorov's work he wrote with G\'acs and Gray, say what the vertical walk is doing in the language of codes. This chapter sets out the algorithmic theory of information as far as the program needs it, proves a coding theorem for the queens that is equivalent to the main equivalence of Chapter~\ref{ch:theorems}, and gives a sufficient condition for derandomizing an $\RP$ algorithm in terms of constructing a random string.

\section{Kolmogorov's three approaches}

Kolmogorov's 1965 paper \cite{Kolmogorov1965}, Three approaches to the quantitative definition of information, distinguishes the combinatorial approach, in which the information in an element of a finite set of size $M$ is $\log M$, the probabilistic approach, which is Shannon's entropy, and the algorithmic approach, which he was introducing. The complexity of a finite object $x$ relative to a universal machine $U$ is
\[
K(x) = \min\{ |p| : U(p) = x \} ,
\]
the length of the shortest program that prints $x$, where $U$ is a universal prefix machine, so that the set of programs on which $U$ halts is prefix free and $K$ is prefix complexity; it differs from plain complexity by $O(\log |x|)$, and the coding theorem and the gambling bound below are statements about the prefix version. The invariance theorem, proved independently by Solomonoff \cite{Solomonoff} in 1964, Kolmogorov in 1965 and Chaitin \cite{Chaitin} in 1966, says that the choice of $U$ changes $K$ by an additive constant, so that the complexity of an object is defined up to a constant that does not depend on the object. The standard reference for everything in this section is Li and Vit\'anyi \cite{LiVitanyi}. Three facts follow at once. $K(x) \le |x| + 2 \log |x| + c$, since a string can be printed from a self delimiting copy of itself. Most strings of length $n$ have $K(x) \ge n - c$, since there are fewer than $2^{n-c}$ programs shorter than $n - c$; these are the incompressible or random strings, and Martin L\"of \cite{MartinLof} showed in 1966 that incompressibility, suitably stated, coincides with passing every effective statistical test. And $K$ is not computable, since a program that computed it could find, for each $n$, the first string of complexity at least $n$, in fewer than $n$ bits for large $n$.

The three approaches agree where they overlap. If $x$ is drawn from a computable distribution $P$ on strings of length $n$ then $\EE_P K(X) = H(P) + O(K(P))$, the expected algorithmic complexity is the Shannon entropy up to the complexity of the distribution, and if $x$ is a typical element of a finite set $S$ then $K(x) \le \log |S| + K(S) + O(1)$ with equality for most elements, the combinatorial measure. Kolmogorov drew the conclusion in his 1983 paper \cite{Kolmogorov1983} on the combinatorial foundations of information theory: probability can be dispensed with in the definition of information, and randomness of an individual object can be defined without a probability space, as the absence of a short description.

Two of the themes of this book begin in Kolmogorov's seminar. Karatsuba\index{Karatsuba, Anatoly}, as a student in 1960, refuted Kolmogorov's conjecture that multiplying two $n$ digit numbers requires $n^2$ operations, with the first subquadratic algorithm, and Kolmogorov wrote it up for him \cite{KaratsubaOfman}; it is the first result of what became algorithmic complexity theory in the Russian school. And Levin\index{Levin, Leonid}, Kolmogorov's student, found in 1973 both the universal search problems of Chapter~\ref{ch:npcomplete} and, in 1974 \cite{Levin1974}, the coding theorem of algorithmic information: the universal a priori probability $\mathbf{m}(x)$, the probability that a universal machine fed random bits prints $x$, satisfies $-\log \mathbf{m}(x) = K(x) + O(1)$, so that the shortest description and the most probable explanation are the same thing. Cover and Thomas call this the algorithmic form of Occam's razor.

\section{Cover on Kolmogorov}

Cover, G\'acs and Gray's 1989 article \cite{CoverGacsGray} in the Annals of Probability, Kolmogorov's contributions to information theory and algorithmic complexity, is the account a probabilist should read. It places the algorithmic theory beside Kolmogorov's work on entropy of dynamical systems and $\epsilon$ entropy, and it states the correspondence between the three approaches with proofs. The fourteenth chapter of Cover and Thomas's Elements of Information Theory \cite{CoverThomas} is the textbook version, and it contains the two results of Cover's own that this chapter uses.

The first is enumerative coding, from Cover's 1973 paper of that name \cite{Cover1973}. Let $S$ be a finite set of strings of length $n$ with a known structure, and index its elements in lexicographic order. The index of $x \in S$ is
\[
i(x) = \sum_{j=1}^{n} \sum_{b < x_j} \big| \{ y \in S : y_1 \dots y_{j-1} = x_1 \dots x_{j-1}, \ y_j = b \} \big| ,
\]
the number of elements of $S$ that precede $x$, computed one coordinate at a time from the number of completions of each prefix. The code is the binary expansion of $i(x)$ in $\lceil \log_2 |S| \rceil$ bits, it is optimal for the uniform distribution on $S$, and encoding and decoding take $n$ evaluations of the prefix completion count for a binary alphabet and up to $n N$ for an alphabet of size $N$. Cover's examples are the strings of fixed weight and the strings with a fixed number of runs, where the counts are binomial coefficients. For the queens the set $S$ is the set of completions of a board, the prefix is a placement of the first $j$ queens, and the completion count of a prefix is the $\sharpP$ complete quantity of Chapter~\ref{ch:completion}. Enumerative coding of the queens is self reducible counting, and the self reducibility\index{self reducibility} tree of Jerrum, Valiant and Vazirani is Cover's code tree.

The second is universal gambling, from Cover's 1974 report \cite{Cover1974} and the chapter of the textbook. A gambler who bets on the successive bits of a sequence $x$, with fair odds and a betting scheme derived from the universal probability $\mathbf{m}$, multiplies the initial wealth by $2^{\,n - K(x) + O(1)}$ after $n$ bits: the log wealth is the compressibility of the sequence. Kelly's growth rate for a known source is the entropy deficit $\log 2 - H$ per bit; the universal gambler attains the algorithmic version without knowing the source. The result is a relative of the universal portfolio of Cover's later work, and in the present setting it says that a sampler for the queens that could be bet against would reveal, by its growth rate, the information the constraints carry. For a uniform solution of the unconditioned problem the compressibility relative to a uniform permutation is $\log N! - \log Q(N) = (\alpha - 1) N + o(N)$ nats, about $1.36 N$ bits, by Simkin's theorem: a solution carries $(\alpha - 1) N$ fewer nats than a permutation, and that is the information the diagonal constraints removed.

\section{The information in a queens solution}

Let $x$ be a solution of the $N$ queens problem, regarded as a permutation. Three descriptions of $x$ are available, with three lengths.

As a permutation, $\log_2 N! \approx N \log_2 N - 1.44 N$ bits, the length of a Lehmer code.

As a solution, $\log_2 Q(N) \approx N \log_2 N - 1.44 \alpha N$ bits, the length of Cover's enumerative code, which is $1.44 (\alpha - 1) N \approx 1.36 N$ bits shorter.

As an output of a short program, $K(x \mid N)$ bits, which for most solutions is $\log_2 Q(N) + O(1)$, since given $N$ the set of solutions has a description of $O(1)$ bits, and for a few solutions, the symmetric ones and the ones found first by a fixed backtracking order, is much less.

The second description can be produced from the first only by counting. This is Cover's code, and the following statement, which we think is new in this form though every part of it is classical, says that the counting is not incidental.

\begin{theorem}[Coding theorem for the queens]
\label{thm:coding}
The following are equivalent.
\begin{enumerate}
\item $\NP = \RP$.
\item There is a constant $c$ and a randomized polynomial time encoder and decoder pair, sharing a random string $\sigma$ of polynomial length, with the following properties on every board $X_0$ with $Q(N; X_0) \ge 1$ completions. With probability at least $1 - 2^{-N^2}$ over $\sigma$, every completion $x$ is encoded by a binary string of a common length $L = L(X_0, \sigma)$, computable in polynomial time from $X_0$ and $\sigma$, with $L \le \log_2 Q(N; X_0) + c \log_2 N$, and decoded correctly. For every $\sigma$, the decoder outputs only completions or a failure symbol, and at most $N^c$ strings of length $L$ decode to any one completion.
\end{enumerate}
\end{theorem}

\begin{proof}
(1) implies (2). If $\NP = \RP$ then by the main equivalence of Chapter~\ref{ch:theorems} there is a fully polynomial randomized approximation scheme for the completion count on every board, and by the median trick its failure probability can be made $2^{-N^3}$ at polynomial cost. Using $\sigma$ as its coins, run it with relative error $1/N$ on the board itself, giving $\tilde N$, and on every prefix $x_1 \dots x_{j-1} b$ for $j \le N$ and $b \le N$, giving $\tilde n_j(b)$; by a union bound over the at most $N^{N+1} + 1$ boards and prefixes that can be queried, each run with the coins $\sigma$, all estimates are within $1 \pm 1/N$ of the truth except with probability $2^{-N^2}$. Encode $x$ by arithmetic coding with the conditional probabilities $p_j(b) = \tilde n_j(b) / \sum_{b'} \tilde n_j(b')$, which is a prefix free code whatever the values of the estimates, with code length $\ell(x) = \sum_j \log_2 (1 / p_j(x_j)) + O(1)$. When the estimates are good, $\ell(x) = \log_2 Q(N; X_0) + O(1)$ for every completion, since the product of the $p_j(x_j)$ telescopes to $1 / Q(N; X_0)$ up to a factor $(1 \pm 1/N)^{2N}$. Set $L = \lceil \log_2 \tilde N \rceil + 4 \log_2 N$ and pad codewords to length $L$ arbitrarily. The decoder recomputes the same estimates from $\sigma$, reads the codeword as the unique prefix of its input that decodes, outputs the completion if $\ell(x) \ge L - 5 \log_2 N$ and fails otherwise; the output is checked to be a completion before it is released. Then at most $2^{L - \ell(x)} \le N^5$ strings decode to $x$ for every $\sigma$, and when the estimates are good every completion satisfies $\ell(x) \ge L - 5 \log_2 N$ and is decoded.

(2) implies (1). Draw $\sigma$ and a uniform string $u$ of length $L$, and decode. A completion $x$ is output with probability at most $2^{-L} N^c + 2^{-N^2}$, the second term from the bad $\sigma$, and at least $(1 - 2^{-N^2}) 2^{-L}$; since $L \le N \log_2 N + c \log_2 N$ the second term is negligible against the first, and the output distribution, conditioned on not failing, is within a factor $2 N^c$ of uniform on the completions, with failure probability at most $1 - N^{-c} / 2$. Fix a good $\sigma$, and take $L$ to be computed in polynomial time from the board and $\sigma$, as in the construction above. Every completion is decoded from some string of length $L$, and distinct completions come from distinct strings, so $Q(N; X_0) \le 2^L$, while the bound on $L$ gives $2^L \le N^c Q(N; X_0)$. Thus $2^L$ approximates the completion count within the polynomial factor $N^c$, and for self reducible problems a polynomial factor approximation can be improved to a fully polynomial randomized approximation scheme, by the argument of Sinclair and Jerrum (1989, Information and Computation 82) \cite{SinclairJerrum} in the framework of Jerrum, Valiant and Vazirani \cite{JVV}. The main equivalence then gives $\NP = \RP$.
\end{proof}

In the language of time bounded complexity, where $K^t(x \mid y)$ is the length of the shortest program that prints $x$ from $y$ within $t(|x|)$ steps, the theorem says that $\NP = \RP$ if and only if every completion of every board has $K^t(x \mid X_0, \sigma) \le \log_2 Q(N; X_0) + O(\log N)$ for a polynomial $t$ and almost every shared string $\sigma$, with the decoder nearly injective; the complexity of a completion given the board is, up to a logarithm, its enumerative code length, if and only if the count can be approximated.

The theorem says that the $(\alpha - 1) N$ nats that separate a permutation from a solution can be saved in polynomial time if and only if $\NP = \RP$. On unconditioned boards, where the count is a function of $N$ alone, the saving is available to anyone who has tabulated $Q(N)$ and the prefix counts, so the theorem is about boards, as every statement in this book is; and the realistic prize of Chapter~\ref{ch:program}, a polynomial Monte Carlo algorithm for $Q(N)$, would give a polynomial enumerative code for the unconditioned solutions as a byproduct, which no one has.

The Markov window computation of Chapter~\ref{ch:open} is the same fact seen through a weaker code. A code whose conditional probabilities for the $j$th queen depend on the previous $k$ queens only is an order $k$ Markov code, and its excess length over Cover's code is the divergence $D_k(N)$ of that chapter, which is linear in $N$ for fixed $k$. The $(\alpha - 1) N$ nats are not stored in any window of the solution. They are stored in the whole of it, which is what a $\sharpP$ complete count should look like from the inside.

\section{Kolmogorov's structure function and the ladder}

Kolmogorov proposed in 1974, in a talk of which only an abstract was published, and whose theory was worked out by Vereshchagin and Vit\'anyi \cite{VereshchaginVitanyi} in 2004, the structure function of an object,
\[
h_x(\kappa) = \min\{ \log_2 |S| : x \in S, \ K(S) \le \kappa \} ,
\]
the log cardinality of the simplest model of complexity at most $\kappa$ that contains $x$. It is a nonincreasing staircase in $\kappa$, it starts at $|x|$ when $\kappa$ is too small for any model but the full cube, and it reaches $K(x) - \kappa$ at the smallest $\kappa$ for which a model of that complexity captures all the structure of $x$; that $\kappa$ is the complexity of the minimal sufficient statistic, and $x$ is a typical element of the corresponding $S$. Vereshchagin and Vit\'anyi proved that every nonincreasing staircase of the right shape is the structure function of some object, and that the minimal sufficient statistic is the model at which the two term code, $K(S)$ bits for the model and $\log |S|$ for the index, first attains $K(x)$.

The level sets of the queens energy are a family of models of the same kind, and the staircase of Chapter~\ref{ch:vertical} is a structure function with the roles of the axes exchanged. The model $\{ S \le m \}$ has complexity $O(\log N)$ for every $m$, and log cardinality $\log N(\le m)$ increasing in $m$; for a solution $x$ the chain of models $\{ S \le m \}$, $m = m_{\max}, \dots, 0$, runs from the full set of permutations down to the solutions, and each step down discards the placements with one more attack. The vertical walk measures the widths of these steps, which are exactly the amounts $\log N(\le m) - \log N(\le m - 1)$ by which the model's log cardinality drops; the sum of the widths is $\log N! - \log Q(N)$, the information in the constraints; and the minimal sufficient statistic of a typical solution, at complexity $O(\log N)$, is the solution set itself, since a typical solution is incompressible given it. The ladder is a structure function whose models are indexed by energy rather than by complexity, and the counting problem is the problem of computing the heights of a structure function whose every model is simple to describe and hard to count.

\section{Randomness, time bounds and P = RP}

The derandomization half of the book has a Kolmogorov form, and it is the cleanest statement we know of a sufficient condition for $\Pc = \RP$.

\begin{proposition}
\label{prop:random}
Let $A$ be an $\RP$ algorithm for a language $L$ using $m = m(n)$ random bits on inputs of length $n$, with acceptance probability at least $\tfrac12$ on $x \in L$, and suppose that for every $x \in L$ the set $B_x$ of rejecting strings has size at most $2^{m/3}$, a strong error bound that we assume rather than derive from amplification. There is a polynomial $t$ such that for every $x \in L$ and every $r \in \{0,1\}^m$ with $\mathrm{CD}^t(r \mid x) \ge 2m/3 + c \log m$, the algorithm accepts: $A(x, r) = 1$. Here $\mathrm{CD}^t(r \mid x)$ is the length of the shortest program that runs in time $t$ given $x$ and accepts $r$ and no other string of length $m$ (distinguishing complexity). Consequently $L \in \Pc$ if there is a deterministic polynomial time procedure that, given $x$, outputs a string $r$ of length $m$ with $\mathrm{CD}^t(r \mid x) \ge 2m/3 + c \log m$ for a suitable polynomial $t$ and constant $c$. Strings of this complexity exist, since all but a fraction $2^{-m/3}$ of the strings of length $m$ have $\mathrm{CD}^t(r \mid x) \ge 2m/3$ by counting programs.
\end{proposition}

\begin{proof}[Proof in outline]
Fix $x \in L$ and let $B_x = \{ r : A(x, r) = 0 \}$ be the rejecting strings, a set of size at most $2^{m/3}$ decidable in polynomial time given $x$. Each $r \in B_x$ can be described, given $x$, by its index in the lexicographic enumeration of $B_x$, which has at most $m/3$ bits, together with a constant size program that enumerates $B_x$ by running $A(x, \cdot)$ on all strings in order and stops at the index; the running time of the description is $2^m \cdot \mathrm{poly}(n)$, which is not polynomial, so this argument gives $K(r \mid x) \le m/3 + O(1)$ but not the time bounded version. For the time bounded version one needs a description that can be checked, not printed, in polynomial time, and this is what the coding arguments descending from Sipser's coding lemma \cite{Sipser} provide. We rely on the bound of Buhrman, Fortnow and Laplante~\cite{BFL}, used in the following form: for a set $B_x$ decidable in polynomial time given $x$, every $r \in B_x$ has $\mathrm{CD}^t(r \mid x) \le 2 \log_2 |B_x| + O(\log m)$ for a suitable polynomial $t$, hence at most $2m/3 + O(\log m)$ here. We do not reprove it, and we note that it bounds distinguishing complexity and does not bound the printing complexity $K^t$, since finding the element of $B_x$ with a given hash value is a search problem not known to be solvable in polynomial time. A string with $\mathrm{CD}^t(r \mid x) \ge 2m/3 + c \log m$ is therefore not in $B_x$. The consequence is immediate: run $A(x, r)$ on the constructed $r$, which accepts exactly when $x \in L$, since on $x \notin L$ the algorithm never accepts.
\end{proof}

The proposition is Sipser's observation of 1983, restated with distinguishing complexity: the rejecting strings of a randomized algorithm are compressible, so an incompressible string is a witness. It is a relative of the hitting set formulation of Chapter~\ref{ch:prp} in Kolmogorov's coordinates, in the regime of strong error bounds, and it explains why that formulation does not become easier when restated. A string of time bounded complexity above $2m/3$ cannot be produced by a short fast program, by definition, and producing one by a short program that runs longer than $t$, which is what the polynomial time procedure in the proposition would be, is possible only if the long program has access to a hard function from which randomness can be extracted, which is the hardness versus randomness trade of Chapter~\ref{ch:pseudo}. The vertical form, that the dense level set of $A$ can be hit by an explicit point, and the Kolmogorov form, that an incompressible string can be written down, are closely related problems, and the second says plainly what kind of object is being asked for: a string with no structure, produced by a structured process.

The modern form of the connection is sharper. Liu and Pass \cite{LiuPass} showed in 2020 that one way functions exist if and only if time bounded Kolmogorov complexity $K^t$ is mildly hard to compute on average, for polynomial $t$; Hirahara and others have since placed the minimum circuit size problem and the complexity of $K^t$ at the centre of the average case theory. In these terms the two halves of this book are two questions about the same function. The Monte Carlo half asks whether a specific $\sharpP$ complete count can be approximated, which by Theorem~\ref{thm:coding} is whether the queens solutions can be compressed to their information content in polynomial time; the derandomization half asks whether a string that cannot be compressed can be constructed in polynomial time. Compression and its impossibility, one for a structured set and one for an unstructured string, are what $\NP = \RP$ and $\Pc = \RP$ come to in Kolmogorov's measure.

\section{What Cover would have measured}

Cover's habit was to ask for the growth rate. For the queens the natural bet is on the next queen: a gambler who sees the first $j$ rows of a uniformly random solution and bets on the column of the $(j+1)$st, at fair odds against the uniform distribution on the $N - j$ free columns, grows wealth at the rate $\log(N - j) - H(X_{j+1} \mid X_1, \dots, X_j)$ per row, and the total log wealth over the board is $\log N! - \log Q(N) = (\alpha - 1) N + o(N)$ nats, the information in the constraints. The gambler who knows only the last $k$ rows grows at a rate smaller by $D_k(N) / N$ per row, which Chapter~\ref{ch:open} computed and found to be a positive constant for every fixed $k$. And a gambler who must compute the odds in polynomial time can grow at the full rate on every board $X_0$, by Theorem~\ref{thm:coding}, if and only if $\NP = \RP$. The growth rate of a polynomial time gambler on the queens is thus an open quantity between $0$ and $\alpha - 1$, and its value is the Monte Carlo half of this book stated as a wager, which is the form in which Cover would have put it.
